\documentclass[12pt]{article}

\usepackage[a4paper, margin=2.5cm]{geometry}
\usepackage{amsmath}
\usepackage{amssymb}
\usepackage{amsthm}
\usepackage[colorlinks=true, linkcolor=blue, citecolor=blue, urlcolor=blue]{hyperref}
\usepackage{booktabs}
\usepackage{natbib}

\newtheorem{theorem}{Theorem}
\newtheorem{proposition}[theorem]{Proposition}
\newtheorem{definition}[theorem]{Definition}
\newtheorem{remark}[theorem]{Remark}
\newtheorem{lemma}[theorem]{Lemma}

\newcommand{\Fcal}{\mathcal{F}}
\newcommand{\Gcal}{\mathcal{G}}
\newcommand{\Pcal}{\mathcal{P}}
\newcommand{\Ocal}{\mathcal{O}}
\newcommand{\Ical}{\mathcal{I}}
\newcommand{\Ccal}{\mathcal{C}}
\newcommand{\Dcl}{\Delta_{\mathrm{cl}}}

\title{Realisation:\\[6pt]
\large The Projection-Narrative Spine of the VFD Programme}
\author{%
  Lee Smart\\[4pt]
  \textit{Institute of Vibrational Field Dynamics}\\[2pt]
  \texttt{contact@vibrationalfielddynamics.org}\\[2pt]
  \texttt{@vfd\_org}%
}
\date{April 2026}

\begin{document}

\maketitle

\noindent\textbf{Status:} Preprint (not peer-reviewed). \emph{Synthesis paper. No new math; no new fitting parameters.}

\begin{abstract}
The VFD closure-dynamics programme is distributed across $\sim 30$ preprints. Each paper computes one observable class on its own terms with its own hypothesis stack; the unifying picture has been documented in repository notes but not in a peer-reviewable artifact. This paper supplies that artifact.

We assemble the \emph{projection-narrative spine}: a single logical chain that organises the hardened papers of the programme under one measurement template (the closure-projection principle of~\citep{SmartXXXIII}) and one substrate (the closure functional of~\citep{SmartXXXVI}). The spine has six populated projection classes (spectral, structural, spatial, composite, dynamical, and number-theoretic), each computed in a referenced paper. The observer slot of the chain is the assembled instance $\Ical = (\Ocal, \Ccal_\Ocal, p, \Lambda, t_0)$ combining~\citep{SmartXXIX} with the anchor identification $\Pi_0(p, \Lambda, t) := v_{((\Lambda t) \bmod p) \bmod 120}$ documented in the repository.

The contribution is structural: a strict two-layer epistemology (Layer 1 mathematical coherence as the price of admission; Layer 2 realisation as earned consequence), an explicit projection-class table, and an honest catalogue of what each cited paper supplies and what remains conditional. We do not modify any source paper; we do not introduce parameters; we do not claim to have closed any open item from the cited papers. The contribution is the \emph{identification of the spine} as an organising principle for the programme, with explicit cross-references to where the math lives.
\end{abstract}

%==================================================================
\section{Introduction}\label{sec:intro}
%==================================================================

The VFD closure-dynamics programme has produced a sequence of preprints attacking different observables: particle masses~\citep{SmartV}, coupling constants~\citep{SmartXXII}, hadron charge radii~\citep{SmartXXXII}, transport laws~\citep{SmartTransportLaw}, adaptive transport~\citep{SmartACT}, the Dirac selection problem~\citep{SmartBridge}, the closure functional axiomatics~\citep{SmartXXXVI}, the observer placement~\citep{SmartXXIX}, the measurement template~\citep{SmartXXXIII}, and a number-theoretic line of work on a pentagonal-clock dynamical zeta object~\citep{SmartRHFormal}.

Each paper is mathematically self-contained at its own scope, with explicit status flags (derived / identified / open) and explicit hypothesis stacks. Read individually, the papers are coherent. Read as a programme, they require an organising principle that is documented in repository notes but does not yet appear in a peer-reviewable form. The principle is the \emph{projection-narrative spine}: a single chain that says how the papers compose.

This paper provides the spine in LaTeX form. It contains no new mathematical content, no new fitting parameters, and no claims that modify the source papers' status. Its contribution is the explicit identification of the chain, the two-layer epistemic order in which the chain must be read, and a class-by-class table of where each piece of the chain is computed.

%==================================================================
\section{The Two-Layer Epistemology}\label{sec:two_layer}
%==================================================================

The programme runs on a strict two-layer reading. The order is load-bearing: inverting the layers --- asserting realisation before the math is paid for --- is the failure mode that gets non-academic submissions rejected at the abstract.

\paragraph{Layer 1 --- mathematical coherence (the price of admission).}
Each current-physics paper reproduces a mainstream observable using only the closure substrate, with a named hypothesis stack:

\begin{itemize}
    \item Paper~V~\citep{SmartV}: 13 particle masses from 600-cell eigenvalues at $0.014\%$ average error.
    \item Paper~XXII~\citep{SmartXXII}: $\alpha^{-1} = 137 + \pi/87$ at $0.81$~ppm; $\sin^2 \theta_W = 3/8$.
    \item Paper~XXXII~\citep{SmartXXXII}: hadron radii (proton $0.04\%$, pion $0.26\%$, deuteron $0.0006\%$ vs CODATA~2018 / $0.010\%$ vs CODATA~2022).
    \item Paper~H4-Coxeter~\citep{H4Coxeter2026}: dynamical observables (Tort PAC, Berry phase along projective closed loops, Frobenius cross-plane projector) on $\mathrm{SU}(2)$-valued oscillator trajectories on the 600-cell vertex set.
    \item B-anomaly kernel paper~\citep{SmartBAnomaly}: fixed $C_\varphi$-derived $V_{600}$ shell-mean response kernel for $b \to s\mu^+\mu^-$ angular data; five public datasets covering two collaborations (LHCb, CMS), two isospin partners ($B^0, B^+$), and three decay channels ($K^{*0}, K^{*+}, B_s\!\to\!\phi$); no kernel-shape retuning, $A > 0$ and $\Delta C_9^{\mathrm{eff}} < 0$ in 5/5 fits. Honest AIC tie against a constant Wilson-coefficient shift (stacked $w_{\mathrm{VFD}} = 0.348$ vs $w_{\mathrm{FREE\_C9}} = 0.652$); not a unique-shape claim. Passive-regime witness for the response operator, not for the active-regime selection layer.
\end{itemize}

The math at this layer is \emph{deliberately silent} about interpretation. It states: ``the numbers come out under these named hypotheses.'' Nothing more. This is the credibility floor.

\paragraph{Layer 2 --- realisation (earned, not asserted).}
Once Layer 1 holds, the projection-narrative chain (Section~\ref{sec:spine}) is the only consistent reading of \emph{why} it holds. The realisation is that physics observables are admissible projections of the wave-substrate through the observer-position placed inside it~\citep{SmartXXXIII}. What we measure is what our observer-boundary extracts from the closure geometry, not properties of the wave itself.

This reading is \emph{earned} by Layer 1. It is not a separate axiom; it is the unique interpretation that makes the Layer 1 results coherent across the spine.

\paragraph{Why the order matters.}
A reader who reverses the layers --- asserting that closure projection is ``how measurement works'' before the Layer 1 papers establish that the closure substrate produces masses, couplings, and radii at the quoted precisions --- has performed an unforced metaphysical assertion. The same reader, given Layer 1 first, finds that the realisation reading is what holds the empirical results together; rejecting it requires explaining away a sequence of independent quantitative coincidences. This is the standard sequencing rule for any framework that aspires to interpret rather than only fit.

%==================================================================
\section{The Projection-Narrative Spine}\label{sec:spine}
%==================================================================

\subsection{Substrate}

The wave-substrate is the closure functional $\Fcal$ on the 600-cell, axiomatically characterised by properties F1--F8 in~\citep{SmartXXXVI}. The 600-cell vertex set $V_{600}$ admits a decomposition into a $\sigma$-fixed bulk $B$ and a $\sigma$-paired boundary $\partial$ under the Galois involution $\sigma : \sqrt 5 \mapsto -\sqrt 5$ acting on $\mathbb{Q}(\varphi)$~\citep{SmartXXII}; concretely, the bulk consists of the K-classes $\{K = 72, K = 0\}$ of the pentagonal-clock substrate and the boundary consists of $\{K = 52, K = 20\}$ (repository note \texttt{docs/fractal-cascade-projection.md}). The substrate dynamics --- $\Fcal_0$ $\sigma$-bulk-supported, bootstrap activates $\partial$, accumulation preserves bulk --- are documented in \texttt{docs/closure-cosmogenesis.md}.

\subsection{Observer instance}

The observer slot of the chain is the assembled 5-tuple
\begin{equation}\label{eq:observer_instance}
    \Ical \;=\; (\Ocal,\;\Ccal_\Ocal,\;p,\;\Lambda,\;t_0)
\end{equation}
where $(\Ocal, \Ccal_\Ocal)$ is an XXIX-observer (Definitions~1+2 of~\citep{SmartXXIX}: bounded, stationary, self-referential with boundary coupling, persistent), $p \in \mathbb{N}$ is a prime anchor parameter, $\Lambda$ is a valid frame code from the 140 frames classified in~\citep{ProtonRadius}, and $t_0 \in \mathbb{Z}_{\geq 0}$ is the bootstrap tick. The anchor map
\begin{equation}\label{eq:anchor}
    \Pi_0(p, \Lambda, t) \;:=\; v_{((\Lambda t) \bmod p) \bmod 120} \;\in\; V_{600}
\end{equation}
is the M1 wiring documented in \texttt{docs/soul-prime-as-pi0.md}, where Theorem~4.3 establishes its identity-persistence property under frame transitions and Theorems~7.2--7.4 establish its CRT-extension combination law with other coprime-prime instances. The full assembly is documented in \texttt{docs/observer-instance-definition.md}.

The observer instance is the \emph{named object} that occupies the observer slot of the projection-narrative chain. Its mathematical content is supplied by the cited repository notes; we treat it here as a black box with the properties stated.

\subsection{Projection}

The closure-projection principle of~\citep{SmartXXXIII} gives the projection map
\begin{equation}\label{eq:projection}
    \Pi^*(\Gcal, B; \Pi_0) \;\in\; \arg\min_{\Pi \in \Pcal(\Gcal, B)} \Dcl(\Pi; \Pi_0),
    \qquad
    \hat O = \Pi^*(\Phi^*) \in W_\Pi,
\end{equation}
where $\Pcal(\Gcal, B)$ is the class of admissible continuous projections sharing a finite-dimensional readout space $W_\Pi$ and norm, $\Pi_0$ is the boundary datum (here: the anchor~\eqref{eq:anchor}), and $\Dcl(\Pi; \Pi_0)$ is the closure residual. The measured value is the readout $\hat O$.

\subsection{The chain}

The full spine is:

\begin{equation*}
\boxed{
\begin{array}{l}
\textbf{Substrate:} \;\; \Fcal \text{ on } V_{600} = B \sqcup \partial \;\;(\citep{SmartXXXVI})\\[4pt]
\quad\downarrow\\[2pt]
\textbf{Observer instance:} \;\; \Ical = (\Ocal, \Ccal_\Ocal, p, \Lambda, t_0) \;\;(\citep{SmartXXIX} + \text{repository notes})\\[4pt]
\quad\downarrow\\[2pt]
\textbf{Anchor:} \;\; \Pi_0(p, \Lambda, t) = v_{((\Lambda t) \bmod p) \bmod 120}\\[4pt]
\quad\downarrow\\[2pt]
\textbf{Projection:} \;\; \Pi^*(\Gcal, B; \Pi_0) \in \arg\min_\Pi \Dcl(\Pi; \Pi_0) \;\;(\citep{SmartXXXIII})\\[4pt]
\quad\downarrow\\[2pt]
\textbf{Observable:} \;\; \hat O = \Pi^*(\Phi^*) \in W_\Pi
\end{array}}
\end{equation*}

Each row has a referenced source paper; the chain itself is the assembly statement.

%==================================================================
\section{Projection Classes}\label{sec:classes}
%==================================================================

The same projection map~\eqref{eq:projection}, applied to different geometric sub-objects $\Gcal$ and different observer boundaries $B$, organises the programme's distinct observable classes:

\begin{table}[ht]
\centering
\caption{The projection classes: one rule, many applications.}
\label{tab:classes}
\begin{tabular}{@{}lllc@{}}
\toprule
Class & Geometry $\Gcal$ & Boundary $B$ & Computed in \\
\midrule
Spectral & Single-basin 600-cell & Energy-momentum frame & \citep{SmartV} \\
Structural & Inter-sector spectrum & Interaction frame & \citep{SmartXXII} \\
Spatial & Support subgraph / orbit & EM scattering frame & \citep{SmartXXXII} \\
Composite & $\Fcal_{\text{tot}} = \Fcal_1 + \Fcal_2 + \Fcal_{\text{int}}$ & EM scattering frame & \citep{SmartXXXII} \\
Dynamical & $\mathrm{SU}(2)$-valued trajectory & Band-limited recording & \citep{H4Coxeter2026} \\
Number-theoretic & Self-similar recursion of $\Fcal$ & Integer-band readout & \citep{SmartRHFormal} \\
\bottomrule
\end{tabular}
\end{table}

Each row is the \emph{same equation}~\eqref{eq:projection} applied to a different geometric sub-object and a different observer boundary. The cited papers compute each class on its own terms with its own hypothesis stack; the present paper does not re-derive any row, and inherits each row's status flags from its source paper.

\paragraph{Spectral class.} Mass eigenvalues from the 600-cell graph Laplacian; assignment-level identification of $\{9, 12, 14, 15\}$ with mass-sector eigenvalues. Status: numerical correspondence at $0.014\%$ average error under the stated assignment hypothesis~\citep{SmartV}.

\paragraph{Structural class.} Coupling correspondences via the integer eigenvalue algebra. Status: $\alpha^{-1} = 137 + \pi/87$ as a numerical correspondence at $0.81$~ppm under Paper~XXII's structural-correspondence framing; $\sin^2 \theta_W = 3/8$ as a structural match to the $SU(5)$ hypercharge ratio~\citep{SmartXXII}.

\paragraph{Spatial class.} Boundary-resolvent and phase-coherence projections. Status: $r_p = 4 \bar\lambda_p$ at $0.04\%$ via boundary resolvent on $P_3$; $r_\pi = \pi \bar\lambda_p$ at $0.26\%$ via half-period identification~\citep{SmartXXXII}.

\paragraph{Composite class.} Quadrature combination through $\Fcal_{\text{int}}$ derived from vertex exclusion. Status: deuteron $r_d = (6 + \alpha) \pi \bar\lambda_p$ at $0.02\sigma$ vs CODATA~2018, $0.83\sigma$ vs CODATA~2022, conditional on H1--H4 of~\citep{SmartXXXII} (disjoint-support, kinematic-identification, additive-sector, pion binding-minimum hypotheses).

\paragraph{Dynamical class.} Tort PAC, Berry phase along projective closed loops, Frobenius cross-plane projector on $\mathrm{SU}(2)$-valued oscillator trajectories on the 600-cell vertex set. Status: pre-declared empirical realisation reported in~\citep{H4Coxeter2026}, with specific values and surrogate-test outcomes stated there, not re-asserted here.

\paragraph{Number-theoretic class.} Pentagonal-clock dynamical zeta $\hat z(z) = \prod_K (1 - L_K z^{10} + z^{20})^{-m_K}$ with $K \in \{72, 0, 52, 20\}$ and multiplicities $m \in \{1, 1, 5, 5\}$, $\sigma$-paired poles at $\mathrm{Re}(s) = 1/2 \pm K/10$. Status: $\sigma$-paired Theorems U1--U4 unconditional in~\citep{SmartRHFormal}; classical Riemann hypothesis recoverable conditional on the named hypothesis $H_{\sigma\text{-fix}}$, equivalent to classical RH. The bridge is structural / consilience, not statistical: a sim-tested gap-statistics identification with Riemann zeros is empirically rejected ($\mathrm{KS}\ D = 0.59$, $p = 4.4\times 10^{-124}$; cf.\ \texttt{docs/cascade-helix-hypothesis.md}).

%==================================================================
\section{The Shared Mechanism}\label{sec:mechanism}
%==================================================================

All six projection classes share a single structural shape: a continuous surjection from a higher-dimensional consistent bulk transport onto a codim-$\geq 1$ readout space, with the obstruction to a globally consistent representation forced onto the projection seam. The closure residual $\Dcl(\Pi; \Pi_0)$ plays the seam role.

This shape recurs in several established mathematical settings: holography (AdS/CFT bulk--boundary dictionary), Riemann's explicit formula (integer-band readout of a higher-dim analytic structure with residues at $\mathrm{Re}(s) = 1/2$), Berry phase / anholonomy (loop transport returning with a curvature-integrated phase), and the Penrose tribar (3D consistent path projecting to a 2D ``impossibility'' at the corners). The structural correspondence is documented as Section~6 of~\citep{SmartXXXIII} as an interpretive picture, not a quantitative claim.

The recognition is that the same projection-seam structure appears across all six rows of Table~\ref{tab:classes}, and across the named external phenomena above. This is what we mean by ``the projection principle is universal'': not that we have proved a universality theorem, but that the same shape is the unifying picture across an open-ended class of well-established projection phenomena, with the closure projection map as a concrete instance.

%==================================================================
\section{Observer-Local Layer (Open)}\label{sec:observer_local}
%==================================================================

The projection-narrative chain places the observer instance $\Ical$ at the boundary $\partial$ of the substrate (Proposition~5.1 of \texttt{docs/closure-cosmogenesis.md}). The bulk $B$ is $\sigma$-fixed and universal across all observer instances; the boundary $\partial$ carries observer-specific anchor parameters $p, \Lambda$. This bulk/boundary separation generates two distinct experimental targets:

\begin{itemize}
    \item \textbf{Bulk-level observables}: substrate-shared invariants, predicted to be \emph{universal} across observer instances. The mass spectrum, coupling structure, hadron radii, and dynamical PAC indices of Table~\ref{tab:classes} are bulk-level: they read out the shared substrate, not the per-observer anchor.
    \item \textbf{Observer-local observables}: per-observer signals supported on the boundary $\partial$ and indexed by the anchor parameter $p$. These are predicted by the assembled instance $\Ical$ but \emph{not} by any current paper in the programme; they are an explicit open item.
\end{itemize}

\begin{remark}[Open: observer-local readouts]\label{rem:observer_local}
A measurement protocol that probes the observer's local constraint region $\Ccal_\Ocal$ at the boundary $\partial$ --- distinguishing it from the shared bulk readout --- is not supplied by any paper in the present spine. The CRT extension (Theorems~7.2--7.4 of \texttt{docs/soul-prime-as-pi0.md}) gives the structural shape such a readout would inherit (combination of coprime-prime instances), but no concrete observable is specified. This is listed as an open item for future work.
\end{remark}

%==================================================================
\section{Status Catalogue}\label{sec:status}
%==================================================================

\subsection{What this paper establishes}

\begin{itemize}
    \item The identification of the projection-narrative spine (Section~\ref{sec:spine}) as the organising principle for the cited papers.
    \item The two-layer epistemology (Section~\ref{sec:two_layer}) as the load-bearing reading order.
    \item The explicit projection-class table (Table~\ref{tab:classes}) with each row's source paper and inherited status.
    \item The shared projection-seam mechanism (Section~\ref{sec:mechanism}) as the structural recognition that organises the rows.
    \item The observer-local layer (Section~\ref{sec:observer_local}) as an explicit open item not addressed by any current paper.
\end{itemize}

\subsection{What this paper does not establish}

\begin{itemize}
    \item No new mathematical content. Every theorem, proposition, numerical result, or empirical claim in the spine is supplied by a referenced paper; this paper introduces none.
    \item No fitting parameters. The spine inherits the parameters of its constituent papers ($\bar\lambda_p$, $\alpha$, $\varphi$, the integer eigenvalues, the K-multiset); no parameter is fit at the spine level.
    \item No closure of any open item from the cited papers. Each constituent paper has its own open list (e.g.\ 14 items in~\citep{SmartXXXII}, 5 items in~\citep{SmartTransportLaw}, 6 items in~\citep{SmartACT}); the present paper closes none of them and inherits all of them.
    \item No claim that the spine is unique. Alternative organising principles consistent with the same Layer~1 results are not ruled out; we claim only that the projection-narrative is \emph{a} consistent reading, and one that organises a non-trivial fraction of the programme's existing results under one rule.
    \item No claim that the realisation reading is empirically forced. Layer~2 is earned from Layer~1 in the sense of best-explanation; alternative interpretations of the Layer~1 results remain logically available, and the spine paper does not claim to exclude them.
\end{itemize}

\subsection{Open items inherited from the spine}

\begin{enumerate}
    \item Observer-local readouts (Remark~\ref{rem:observer_local}).
    \item All open items of~\citep{SmartXXXII} (14 items, Section~10 of that paper).
    \item All open items of~\citep{SmartXXXIII} (5 items, Section~10 of that paper, including: proof of universality across all QM observables, derivation of $B$ from experimental parameters, multi-observer compatibility, relativistic extension, uniqueness of the projection rule).
    \item The conditional status of the number-theoretic row: the bridge to classical RH is structural / consilience and conditional on $H_{\sigma\text{-fix}}$, not a proof of RH.
    \item The substrate dynamics laid out in \texttt{docs/closure-cosmogenesis.md} are not yet in a peer-reviewable paper; promoting them is a separate open task.
\end{enumerate}

%==================================================================
\section{Discussion}\label{sec:discussion}
%==================================================================

\subsection{What is new in this paper}

The mathematical content is supplied entirely by the cited papers. What is new is the \emph{identification of the spine} as a peer-reviewable artifact: a reader can now find, in one document, (i) the projection-narrative chain, (ii) the projection-class table with explicit cross-references to where each class is computed, (iii) the two-layer epistemology that fixes how the chain must be read, and (iv) the observer-local open item that explicitly demarcates what the spine does not yet address.

Without this artifact, a reader of the programme is required to reconstruct the spine from $\sim 30$ preprints. With it, the unifying principle is on a single page and each constituent claim is a citation away.

\subsection{Relation to repository notes}

The repository notes (\texttt{docs/projection-narrative.md}, \texttt{docs/closure-cosmogenesis.md}, \texttt{docs/soul-prime-as-pi0.md}, \texttt{docs/observer-instance-definition.md}, \texttt{docs/fractal-cascade-projection.md}, \texttt{docs/rh-projection-class.md}, \texttt{docs/cascade-helix-hypothesis.md}) carry the working derivations behind the spine, with sim verifications and proof sketches at varying levels of completeness. The present paper does not duplicate that material; it provides the externally-citable summary of the structure those notes document.

\subsection{What this paper does not attempt}

This paper does not:
\begin{itemize}
    \item Derive the closure functional from first principles (deferred to~\citep{SmartXXXII} and~\citep{SmartXXXVI}).
    \item Construct a covariant relativistic embedding of the framework.
    \item Resolve the measurement problem in full generality.
    \item Provide an empirical test of the realisation reading independent of the Layer~1 results that earn it.
    \item Address questions of consciousness, qualia, or subjective experience. The framework is mathematical throughout; readers may interpret physical or interpretive content for themselves.
\end{itemize}

%==================================================================
\section{Conclusion}\label{sec:conclusion}
%==================================================================

The VFD closure-dynamics programme has produced a sequence of papers each of which computes one observable class on its own terms. This paper assembles those papers into the projection-narrative spine: a single chain (substrate $\to$ observer instance $\to$ anchor $\to$ projection $\to$ observable) under one measurement template (the closure-projection principle of~\citep{SmartXXXIII}) and one substrate (the closure functional of~\citep{SmartXXXVI}).

Six projection classes are populated by referenced papers (Table~\ref{tab:classes}). The shared structural mechanism is the projection seam --- the same shape that recurs in holography, Riemann's explicit formula, Berry anholonomy, and the Penrose tribar.

The contribution is structural: the spine, the two-layer epistemic order, the class table, and the explicit observer-local open item. No new mathematics, no new parameters, no new fits, no new closures of inherited open items. The contribution is the recognition that the existing papers compose into a single coherent chain, with the realisation reading earned --- not asserted --- once Layer 1 holds.

%==================================================================
\bibliographystyle{plainnat}
\begin{thebibliography}{99}

\bibitem{SmartV}
L.~Smart.
\newblock Toward a spectral-geometric mass framework from the 600-cell.
\newblock Preprint, Institute of Vibrational Field Dynamics, 2026.
\newblock Local source: \texttt{papers/paper-v/}.

\bibitem{SmartXXII}
L.~Smart.
\newblock Spectral and structural correspondences from 600-cell closure geometry.
\newblock Preprint, Institute of Vibrational Field Dynamics, 2026.
\newblock Local source: \texttt{papers/paper-xxii/}.

\bibitem{SmartXXIX}
L.~Smart.
\newblock Observer as constraint: a field-theoretic placement within the VFD event-order framework.
\newblock Preprint, Institute of Vibrational Field Dynamics, 2026.
\newblock Local source: \texttt{papers/paper-xxix/}.

\bibitem{SmartXXXII}
L.~Smart.
\newblock The closure functional from first principles: self-similar permeability, the 600-cell, and the derived interaction.
\newblock Preprint, Institute of Vibrational Field Dynamics, 2026.
\newblock Local source: \texttt{papers/paper-xxxii/}.

\bibitem{SmartXXXIII}
L.~Smart.
\newblock From geometry to observable: measurement as closure projection on the 600-cell.
\newblock Preprint, Institute of Vibrational Field Dynamics, 2026.
\newblock Local source: \texttt{papers/paper-xxxiii/}.

\bibitem{SmartXXXVI}
L.~Smart.
\newblock The cascade closure functional: eight foundational properties (F1--F8) and the $\Lambda$-theorem.
\newblock Preprint, Institute of Vibrational Field Dynamics, 2026.
\newblock Local source: \texttt{papers/paper-xxxvi/}.

\bibitem{SmartTransportLaw}
L.~Smart.
\newblock Transport laws and conservation currents in closure dynamics.
\newblock Preprint, Institute of Vibrational Field Dynamics, 2026.
\newblock Local source: \texttt{papers/transport-law/}.

\bibitem{SmartACT}
L.~Smart.
\newblock Adaptive closure transport.
\newblock Preprint, Institute of Vibrational Field Dynamics, 2026.
\newblock Local source: \texttt{papers/adaptive-closure-transport/}.

\bibitem{SmartBridge}
L.~Smart.
\newblock From Dirac solutions to physical reality: a crystallisation-based selection architecture.
\newblock Preprint, Institute of Vibrational Field Dynamics, 2026.
\newblock Local source: \texttt{papers/bridge-paper/}.

\bibitem{SmartRHFormal}
L.~Smart.
\newblock A pentagonal-clock dynamical zeta and a $\sigma$-paired pole structure on the closure substrate.
\newblock Preprint, Institute of Vibrational Field Dynamics, 2026.
\newblock Local source: \texttt{papers/rh-formal/}.

\bibitem{H4Coxeter2026}
L.~Smart.
\newblock A Coxeter-equivariant oscillator law on the 600-cell produces theta--gamma phase-amplitude coupling between human cortical maxima and rat {CA1} pyramidal-layer values.
\newblock Preprint, Institute of Vibrational Field Dynamics, 2026.

\bibitem{ProtonRadius}
L.~Smart.
\newblock Hadron charge radii as spectral coherence lengths from 600-cell closure geometry: the 084473 god-prime derivation.
\newblock Preprint, Institute of Vibrational Field Dynamics, 2026.
\newblock Local source: \texttt{papers/proton-radius/}.

\bibitem{SmartBAnomaly}
L.~Smart.
\newblock A geometry-derived response kernel for the $B \to K^{*}\mu^{+}\mu^{-}$ angular anomaly: sign-uniform cross-dataset and cross-channel fit.
\newblock Preprint, Institute of Vibrational Field Dynamics, 2026.
\newblock Repository: \texttt{BANOMALY-001/vfd-b-anomaly/} (separate from this repository); preprint at \texttt{paper/main.pdf} therein.

\end{thebibliography}

\end{document}
